\usepackage{xspace}
%\newcommand{\OpenLB}[1][!]{\textit{\sffamily{\ifx!#1OpenLB\else OpenLB~#1\fi}}\xspace}

\usepackage[utf8]{inputenc} %% utf8 support

\usepackage[dvipsnames]{xcolor}

\usepackage{scrhack} %% kill toc warnings

%% helvet (phv) and palatino (ppl) only declare T1/OT1 font shapes. Under XeLaTeX
%% fontspec switches the encoding to TU, which left every one of their shapes
%% undefined -- see "LaTeX Font Warning: Font shape `TU/ppl/b/n' undefined". The
%% substitution dropped the bold series, so \textbf produced upright text
%% throughout the PDF. Load them only for the engine that can use them.
\usepackage{iftex}
\ifPDFTeX
  \usepackage{helvet} %% modern font (for headings)
  \usepackage{palatino} %% palatino font (for text)
\fi

\setkomafont{disposition}{\normalfont\sffamily\normalcolor\bfseries} %% regulate spacing of headings

\usepackage[onehalfspacing]{setspace} %% linespacing

\usepackage{amsmath,amsfonts,amsthm}
\theoremstyle{plain}
\newtheorem{theorem}{Theorem}
\usepackage{booktabs}


\usepackage{geometry}
\usepackage{tabularx}

\usepackage{graphicx}
\usepackage{url}
\usepackage{hyperref}
\hypersetup{
    colorlinks,
    linkcolor={red!50!black},
    citecolor={blue!50!black},
    urlcolor={blue!80!black}
}

\usepackage{listings}
\usepackage{algorithm,algpseudocode}
\usepackage{longtable}
\usepackage{pdflscape}
\usepackage{color}
\usepackage{colortbl}

\usepackage{placeins}
\usepackage{xspace}
\usepackage{ifthen}

\usepackage{pdfpages}

%\renewcommand{\abstractnamefont}{\normalfont\sffamily\Large\bfseries}

\definecolor{gray}{rgb}{0.4,0.4,0.4}
\definecolor{darkblue}{rgb}{0.0,0.0,0.6}
\definecolor{cyan}{rgb}{0.0,0.6,0.6}

\definecolor{lighttuerkis}{HTML}{F7FCFC}
\definecolor{tuerkis}{HTML}{D9EEEE}
\definecolor{RoyalBlue}{RGB}{22,111,191} %%166FBF {0.086,0.435,0.749}
\definecolor{RoyalBluedark}{RGB}{10,80,191}
\definecolor{lightgray}{RGB}{240,240,240}
\definecolor{darkgray}{RGB}{225,225,225}

\definecolor{lightantituerkis}{HTML}{FFFCFA}
\definecolor{antituerkis}{HTML}{FFE5CF}
\definecolor{darkantituerkis}{HTML}{FF9A0C} %%{1,.604,.047}

\usepackage{tikz}
\usepackage{subcaption}
\usepackage{pgfplots}
\pgfplotsset{compat=newest}
\usetikzlibrary{
    automata,
    arrows.meta,
    calc,
    shapes.geometric,
    graphs,
    positioning,
    calligraphy,
    decorations.pathreplacing,
    angles,
    external,
    quotes
}
%\tikzexternalize[prefix=./figuresOut/]
\tikzstyle{startstop} = [ellipse, minimum width=3cm, minimum height=1cm, text centered, draw=black, fill=blue!20]
\tikzstyle{process} = [rectangle, minimum width=2cm, minimum height=1cm, text centered, draw=black, fill=gray!20]
\tikzstyle{decision} = [diamond, minimum width=1cm, minimum height=1cm, text centered, draw=black, fill=red!20]
\tikzstyle{arrow} = [thick,->,>=stealth]


\usepackage[
  table-parse-only,
  detect-all,
  list-separator = \text{;\ },
  list-units = single,
  range-units = single
]{siunitx}
\usepackage{unicode-math}

% \lstdefinelanguage{XML} %% enables syntax highlithning of xml code
% {
%   numbers=none,
%   captionpos=b, %% label position t-top or b-bottom
%   basicstyle=\footnotesize\singlespacing,
%   showstringspaces=false,
%   morestring=[b]",
%   morestring=[s]{>}{<},
%   morecomment=[s]{<?}{?>},
%   stringstyle=\color{black},
%   identifierstyle=\color{darkblue},
%   keywordstyle=\color{cyan},
%   morekeywords={timestep, index, group, part, byte_order, file ,version, type}% list your attributes here
% }

% \lstdefinelanguage{bash} %% enables syntax highlithning of bash code
% {
%   numbers=none,
%   captionpos=b, %% label position t-top or b-bottom
%   basicstyle=\small\ttfamily,
%   showstringspaces=false,
%   keywordstyle=\color{RoyalBlue}\bfseries, % Add \bfseries for bold
%   morekeywords={convert,ffmpeg}
% }

% \lstdefinelanguage{myc++} %% enables syntax highlighting of c++ code
% {
%   language=C++,
%   captionpos=b, %% label position t-top or b-bottom
% %  frame=single,
%   keywordstyle=\color{RoyalBlue}\bfseries\em, % Add \bfseries for bold
%   keywordstyle=[2]\color{RoyalBlue}\bfseries\em,
%   basicstyle=\footnotesize\ttfamily,
% %  basicstyle=\small\ttfamily, % Makes listings a smaller font size and a different 
% %  identifierstyle=\color{NavyBlue}, % Color of text inside brackets
%   commentstyle=\color{gray}\ttfamily, % Color of comments
%   stringstyle=\rmfamily, % Font type to use for strings
%   numbers=left, % Change left to none to remove line numbers
%   numberstyle=\scriptsize, % Font size of the line numbers
%   stepnumber=1, % Increment of line numbers
%   numbersep=8pt, % Distance of line numbers from code listing
%   showstringspaces=false, % Sets whether spaces in strings should appear underlined
%   breaklines=true, % Force the code to stay in the confines of the listing box
%   lineskip=1pt, %% linespacing
% %  frameround=ftff, % Uncomment for rounded frame
%   frame=tB, % Frame border - none/leftline/topline/bottomline/lines/single/shadowbox/L
%   belowcaptionskip=.75\baselineskip, % Space after the "Listing #: Desciption" text and the listing box
% %  keywords={void, if, endif, else, for, endfor, foreach, while, endwhile, class, Class, bool, const, double, auto, T, S, begin, end, int, char, std::string, std::shared_ptr, DESCRIPTOR, PARTICLETYPE, virtual, unsigned, ADDESCRIPTOR, NSDESCRIPTOR, PARTICLE, Vector, constexpr, MOMENTA, EQUILIBRIUM, PARAMETERS, CELL},
% %  morekeywords=[2]{CuboidStructureXD, CuboidDecomposition, SuperCommunicator, ConcreteBlockLattice, FieldArrayD, SuperParticleSystemXD, Data, Timer, SuperParticleSystem3D, ParticleSystem3D, ParticleSystemXD, BlockLatticeXD, BlockLattice3D, BlockLattice2D, BlockLattice, Cell, SuperLatticeXD, SuperLattice, SuperLattice, SuperLattice, Triangle, XMLreader, STLreader, Octree, Force3D, Boundary3D, IndicatorF3D, Particle3D, ContactDetection, MpiNonBlockingHelper, SuperParticleSystem3D, SuperParticleSysVtuWriter, SuperLatticeInterpPhysVelocity3D, StokesDragForce3D, MaterialBoundary3D, BGKdynamics, LBconverter, Serializable, Dynamics, BlockStructureXD, IndicatorFXD, IndicatorCuboid3D, IndicatorCuboid2D, IndicatorIdentity2D, SuperParticleSystem, LoadBalancer, HeuristicLoadBalancer, AdvectionDiffusionBGKdynamics, HeuristicLoadBalancer, SuperVTMwriter3D, SuperLatticePhysVelocity3D, SuperLatticeVelocity3D, SuperLatticeDensity3D, SuperLatticeFfromAnalyticalF3D, SuperLatticePhysPressure3D, SuperGeometryF3D, SuperLatticeCuboid3D, SuperL2Norm3D, SuperEuklidNorm3D, SuperEuklidNorm2D, BlockReduction3D2D, BlockReduction2D2D, Hyperplane3D, HyperplaneLattice3D, BlockGifWriter, SuperPlaneIntegralF3D, SuperPlaneIntegralFluxVelocity3D, SuperPlaneIntegralFluxPressure3D, SuperPlaneIntegralFluxF3D, SuperPlaneIntegralFluxMass3D, UnitConverter, UnitConverterFromResolutionAndLatticeVelocity, OstreamManager, GenericF, AnalyticalF, AnalyticalF3D, AnalyticalF2D, AnalyticalConst2D, AnalyticalConst3D, AnalyticPlus2D, AnalyticalIdentity2D, SmoothIndicatorF2D, SmoothIndicatorF3D, IndicatorSphere3D, IndicatorCircle3D, Indicator3D, CuboidDecomposition3D, SuperGeometryXD, SuperGeometry3D, SuperGeometry2D, SuperGeometry, SuperLatticeRank3D, PolynomialStartScale, SquarePoiseuilleInflow3D, SuperLatticeF3D, SuperLatticeFlux3D, SuperLatticePhysVelocityFlux3D, SuperLatticePhysPressureFlux3D, SuperExternal3D, Poiseuille2D, SuperIndicatorMaterial3D, SuperIndicatorFXD, SuperRelativeErrorLpNormXD, SuperAbsoluteErrorLpNormXD, SuperLatticeKnudsen2D, SuperLatticeRefinementMetricKnudsen2D, LATTICE_DESCRIPTOR, PARTICLE_DESCRIPTOR, FIELD_BASE, D2Q9, D3Q19, POPULATION, FORCE, POROSITY, VELOCITY_NUMERATOR, VELOCITY_DENOMINATOR, GENERAL, REVISION, OMEGA, ResolvedParticle3D, ResolvedSphere3D, ParticleSystem, ParticleManager, VerletParticleDynamics,DynamicsTupleParser,OSMParser}
% }

% \lstdefinestyle{floater}{
%     float=tb,
%     frame=tB,
%     language=myc++
% }

% \lstdefinestyle{intext}{
%     frame=leftline,
%     numbers=none,
%     xleftmargin=10pt,
%     language=myc++
% }

% \usepackage[T1]{fontenc}

% % \usepackage[
% %   backend       = biber,
% %   citetracker   = true,
% %   defernumbers  = true,
% % %  sorting      = ydnt,
% %   sorting      = nyt,
% %   natbib        = true,
% %   style         = numeric-comp,    %% authoryear or alphabetic
% %   giveninits    = true,
% %   maxnames      = 2,
% %   minnames      = 1,
% %   maxbibnames   = 14,
% %   url           = false,
% %   isbn          = false,
% %   doi           = true,
% %   eprint        = false
% %   ]
% %   {biblatex}

% %\addbibresource[datatype=bibtex]{olb_ug_mathias.bib}
% %\addbibresource[datatype=bibtex]{olb_ug_external.bib}


% %% cut url and texttt to pagewidth
% \apptocmd{\UrlBreaks}{\do\f\do\m}{}{}
% \setcounter{biburllcpenalty}{7000} %% Kleinbuchstaben
% \setcounter{biburlucpenalty}{7000} %% Großbuchstaben
% \setcounter{biburlnumpenalty}{7000} %% Zahlen

% \newcommand*\justify{
%   \fontdimen2\font=0.4em %% interword space
%   \fontdimen3\font=0.2em %% interword stretch
%   \fontdimen4\font=0.1em %% interword shrink
%   \fontdimen7\font=0.1em %% extra space
%   \hyphenchar\font=`\- %% allowing hyphenation
% }

\newcommand{\lb}{$\left<\right.$}
\newcommand{\rb}{$\left.\right>$}
%% Math macros shared with the website live in macros.tex.
\newcommand{\class}[1]{\sloppy\lstinline[language=myc++]{#1}}
\newcommand{\struct}[1]{\sloppy\lstinline[language=myc++]{#1}}
%% bm cannot embolden the math characters of unicode-math (loaded above), so
%% \bm{\Pi} stops XeLaTeX with "Improper alphabetic constant". Use the bold
%% italic alphabet of unicode-math instead.
\renewcommand{\boldsymbol}[1]{\symbfit{#1}}

\let\OldTexttt\texttt % make \texttt sloppy to respect page margins
\renewcommand{\texttt}[1]{\sloppy{\footnotesize\OldTexttt{#1}}}

\newcommand{\eg}{e.g.\ }
\newcommand{\ie}{i.e.\ }
\newcommand{\TODO}[1]{\textcolor{red}{[\textbf{#1}]}}
%\newcommand{\OpenLB}[1][\@empty]{\textit{\sffamily{\ifthenelse{\isempty{#1}}{OpenLB}{OpenLB~#1}}}\xspace}
\newcommand{\OpenLB}[1][!]{\textit{\sffamily{\ifx!#1OpenLB\else OpenLB~#1\fi}}\xspace}

%\usepackage{tocloft} %% control spacing in table of contents
%\usepackage{mathtools}

%for quarto algorithm
\renewcommand{\Return}{\State \textbf{return}~}
\newcommand{\Print}{\State \textbf{print}~}
\newcommand{\Break}{\State \textbf{break}}
\newcommand{\Continue}{\State \textbf{continue}}
\newcommand{\True}{\textbf{true}}
\newcommand{\False}{\textbf{false}}
\renewcommand{\And}{\textbf{and}~}
\newcommand{\Or}{\textbf{or}~}
\newcommand{\Not}{\textbf{not}~}
\newcommand{\To}{\textbf{to}~}
\newcommand{\DownTo}{\textbf{downto}~}

\usepackage{tocloft} %% control spacing in table of contents

\setcounter{secnumdepth}{3} %% set numbering of sections:
\setcounter{tocdepth}{2} %% set appearance in table of contents
%% -1 part
%% 0 chapter
%% 1 section
%% 2 subsection
%% 3 subsubsection
%% 4 paragraph
%% 5 subparagraph

\usepackage{mathtools}